\chapter*{Abstract}
\addcontentsline{toc}{chapter}{Abstract}
\begin{otherlanguage}{english}

This capstone project addresses the design and development of an artificial-intelligence-based early warning system for road traffic crashes, carried out under the agreement between the Software Research and Development and Artificial Intelligence Laboratory (LIDeSIA) of the Faculty of Exact, Physical and Natural Sciences of the National University of Córdoba (FCEFyN-UNC) and the Public Prosecutor's Office of the Province of Córdoba. The problem stems from the need of the Road Crash Judicial Unit to allocate specialized physical and human resources (judicial officers, forensic experts, police and health personnel) to events whose occurrence is very hard to anticipate, which forces the institution to oversize its operation.

The project was managed with a hybrid agile methodology: a discovery stage built on weekly Scrum-style iterations with the client, followed by an execution stage organized on a Kanban board. Starting from 4,320 geo-referenced crashes recorded between 2022 and 2024, the dataset was restructured from a fact-table format into a supervised autoregressive formulation, aggregating the information into discrete six-hour blocks and sliding windows. Several artificial intelligence models were tested (LightGBM, multilayer perceptron, geo-referenced models and the automated AutoGluon-TimeSeries pipeline), and an XGBoost model with Poisson regression followed by a decision layer was adopted as the baseline, prioritizing explainability, in line with ISO/IEC 42001.

In the test set, 71.1\,\% of the blocks recorded at least one victim, so the institution's current situation is equivalent to a permanent alert. Against that reference, the model recognized 38.5\,\% of the blocks without victims while keeping a sensitivity of 87.4\,\% on critical blocks, and reduced the squared error by 6\,\% relative to a constant prediction. This is a modest predictive signal. Its main value is to establish a measurable baseline and to validate the feasibility of the approach, rather than to serve yet as a planning tool. Feature importance analysis showed that risk is mainly explained by the recent behavior of the series and the time slot, while weather variables proved marginal. The work includes risk management and an impact assessment in accordance with ISO/IEC 42001 and ISO/IEC 23894, reaching technology readiness level TRL~3 (ISO 16290:2013), and proposes a ten-month implementation plan towards TRL~7, including its schedule, the required profiles and indicators to assess its impact.

\medskip
\noindent\textbf{Keywords:} early warning, road traffic crashes, machine learning, time series, XGBoost, ISO/IEC 42001, product management, agile methodologies.

\end{otherlanguage}
